% Generated draft: requires mathematical and source review.
\documentclass[11pt]{article}
\usepackage[margin=25mm]{geometry}
\usepackage{fontspec}
\setmainfont{Latin Modern Roman}
\usepackage{amsmath,amssymb,mathtools,unicode-math}
\setmathfont{Latin Modern Math}
\usepackage{enumitem,needspace,xurl,hyperref}
\setlength{\parindent}{0pt}
\setlength{\parskip}{5pt}
\setlength{\emergencystretch}{3em}
\newcommand{\N}{\mathbb{N}}
\newcommand{\Z}{\mathbb{Z}}
\newcommand{\Q}{\mathbb{Q}}
\newcommand{\R}{\mathbb{R}}
\newcommand{\C}{\mathbb{C}}
\newcommand{\F}{\mathbb{F}}
\newcommand{\A}{\mathbb{A}}
\newcommand{\PP}{\mathbb{P}}
\newcommand{\E}{\mathbb{E}}
\newcommand{\eps}{\varepsilon}
\newcommand{\dd}{\,\mathrm{d}}
\begin{document}
\section*{UnsolvedMath ID: 701691}
\textbf{Source problem code:} catalogue entry
\subsection{Definitions and mathematical statement}
Fix an oracle set \(X\subseteq\N\). A finitely generated group \(H\) has an \emph{\(X\)-computable presentation}, in the terminology of the supplied record, if there exist a finite set \(A\) and a set \(R\) of words in the free group \(F(A)\), computably enumerable by a Turing machine with oracle \(X\), such that
\[
H\cong F(A)/\langle\!\langle R\rangle\!\rangle_{F(A)}.
\]
Here \(\langle\!\langle R\rangle\!\rangle_{F(A)}\) denotes the normal closure of \(R\) in \(F(A)\). Thus the terminology requires oracle-relative computable enumerability of defining relations, without requiring a decidable word problem.

The question is whether, for every \(X\subseteq\N\), there exists a finitely generated group \(G_X\) with an \(X\)-computable presentation such that, for every finitely generated group \(H\),
\[
\begin{aligned}
&H\text{ has an }X\text{-computable presentation}\\
&\quad\Longleftrightarrow\quad
\exists n\geq 1\;\exists S\subseteq G_X^{*n},\quad
|S|<\infty\ \text{ and }\
H\hookrightarrow G_X^{*n}/\langle\!\langle S\rangle\!\rangle_{G_X^{*n}}.
\end{aligned}
\]
Here \(G_X^{*n}\) is the free product of \(n\) copies of \(G_X\), and \(\hookrightarrow\) means that an injective group homomorphism exists. Its image is a finitely generated subgroup because \(H\) is finitely generated. The integer \(n\), the finite set \(S\), and the embedding may depend on \(X\) and \(H\); no uniform computability requirement for these choices is stated in the source.

As background, the supplied record states that Higman's Embedding Theorem permits the choice \(G_X=\Z\) when \(X\) is computable. This is retained as supplied source context. The attribution is unresolved: the registry names J. Grochow as proposer, whereas the statement begins with the parenthetical attribution ``(A. Bauer)''.
\subsection{Short English statement}
For each oracle \(X\subseteq\N\), does there exist a single finitely generated group \(G_X\), presented using an \(X\)-computably enumerable set of relations, such that the finitely generated groups with such presentations are exactly those embedding into quotients of finite free products of copies of \(G_X\) by finitely normally generated subgroups? The supplied attribution to J. Grochow and the parenthetical reference to A. Bauer require clarification.
\subsection{Sources}
\begin{itemize}
\item Local UnsolvedMath registry, permanent ID \(701691\), source problem code \texttt{catalogue entry}, title \emph{Kourovka Notebook Problem 20.33}. The supplied record labels the problem open; that label is not an independently verified literature assessment.
\item \emph{The Kourovka Notebook (Unsolved Problems in Group Theory)}, 20th issue (2022), problem 20.33; the supplied background also cites the 21st issue edition (2026), arXiv:\texttt{1401.0300v48}. The supplied collection description names E. I. Khukhro and V. D. Mazurov as editors. These bibliographic details have not been independently verified here.
\item Attribution as supplied: J. Grochow, proposed in 2022, in the registry metadata and background; ``(A. Bauer)'' at the start of the statement. The respective roles are unspecified.
\item The supplied statement invokes Higman's Embedding Theorem for the computable-oracle case, but supplies no separate bibliographic reference for that theorem.
\end{itemize}
\end{document}
